\section{Chinese remainders become Lagrange interpolation}
\label{sec:crt-interpolation}

CRT builds an integer with prescribed remainders. Interpolation builds
a polynomial with prescribed values. Polynomial division makes these
the same kind of problem.

Let $K$ be a field, such as $\Q$, $\R$, or $\F_p$. Write $K[X]$ for
polynomials in an indeterminate $X$ with coefficients in $K$. Division
by $X-c$ gives
\[
 f(X)=(X-c)q(X)+r,
\]
where $r$ is constant. Substituting $X=c$ shows $r=f(c)$. In other words,
\[
 f(X)\congr b\pmod{X-c}\quad\Longleftrightarrow\quad f(c)=b.
\]
Here polynomial congruence means that the difference is divisible by the
polynomial modulus.

\subsection*{The same construction, with polynomial moduli}

The division algorithm works in $K[X]$: cancel the leading term using
division in $K$, subtract, and continue while the degree decreases.
Extended Euclid therefore works too. For coprime polynomials $A,B$,
it produces $U,V$ with $UA+VB=1$.

Consequently, for pairwise coprime monic moduli $m_1,\ldots,m_r$, set
$P=\prod_i m_i$ and $P_i=P/m_i$, and use extended Euclid to find
$U_iP_i+V_im_i=1$. The polynomial
\[
 f=\sum_i b_iP_iU_i
\]
has prescribed polynomial remainders $b_i$ modulo $m_i$.
The proof is the integer CRT proof: reduce at each modulus separately.
Uniqueness is modulo $P$, since the pairwise coprime factors $m_i$ all
divide the difference of two solutions. Dividing by the monic polynomial
$P$ selects one representative of degree less than $\deg P$.

For interpolation, take $m_i=X-c_i$ with the $c_i$ distinct. Two such
linear polynomials are coprime: subtracting them gives the nonzero
constant $c_j-c_i$, which is invertible in $K$.
In this special case, the inverse required by CRT is particularly simple.
For
\[
 P_i(X)=\prod_{j\ne i}(X-c_j),
\]
reduction modulo $X-c_i$ gives the nonzero number
$P_i(c_i)=\prod_{j\ne i}(c_i-c_j)$.
Thus define
\[
 \ell_i(X)=\frac{P_i(X)}{P_i(c_i)}
 =\prod_{j\ne i}\frac{X-c_j}{c_i-c_j}.
\]
These are polynomial selectors: $\ell_i(c_i)=1$ and
$\ell_i(c_j)=0$ for $j\ne i$.

\begin{theorem}[Lagrange interpolation]
Let $c_1,\ldots,c_r$ be distinct points of a field $K$.
Given values $b_i\in K$, the unique polynomial of degree less than $r$
with $f(c_i)=b_i$ for every $i$ is
\[
 f(X)=\sum_{i=1}^r b_i\ell_i(X).
\]
Every polynomial with these values has the form
$f(X)+h(X)\prod_i(X-c_i)$, with $h\in K[X]$.
\end{theorem}

\begin{proof}
The selector values prove existence. If $f,g$ have the required values,
each $X-c_i$ divides $f-g$, and hence so does their product. A nonzero
multiple of that degree-$r$ product cannot have degree less than $r$.
This proves uniqueness and the description of all solutions.
The same argument proves the useful \emph{root bound}: a nonzero
polynomial of degree $d$ over a field has at most $d$ distinct roots.
\end{proof}

\begin{center}
\small
\begin{tabularx}{\textwidth}{p{.18\textwidth}XX}
\toprule
 & Integer CRT & Lagrange interpolation \\
\midrule
Local data & $x\congr a_i\pmod{n_i}$ & $f(c_i)=b_i$ \\
Moduli & pairwise coprime $n_i$ & distinct linear factors $X-c_i$ \\
Selector & $N_i(N_i^{-1}\bmod n_i)$ & $P_i(X)/P_i(c_i)$ \\
Formula & $\sum_i a_i e_i$ & $\sum_i b_i\ell_i(X)$ \\
Ambiguity & multiples of $N$ & multiples of $\prod_i(X-c_i)$ \\
Choice & $0\le x<N$ & $\deg f<r$ \\
\bottomrule
\end{tabularx}
\end{center}

\begin{example}[Three values determine a quadratic]
Over $\Q$, prescribe $f(0)=1$, $f(1)=2$, $f(2)=5$. The selectors are
\[
 \ell_0=\frac{(X-1)(X-2)}2,\qquad
 \ell_1=-X(X-2),\qquad
 \ell_2=\frac{X(X-1)}2.
\]
Combining them gives $f=\ell_0+2\ell_1+5\ell_2=X^2+1$.
All other solutions are $X^2+1+X(X-1)(X-2)h(X)$.

The same computation works in $\F_5$, where the prescribed values are
$1,2,0$ and division by $2$ means multiplication by $3$.
The certificate $2\cdot3-5=1$ makes that division legitimate.
The result is again $X^2+1$, now with coefficients modulo $5$.
\end{example}

\begin{warningbox}
Distinct interpolation points must have invertible differences.
This is automatic in a field. It can fail for composite moduli:
over $\Z/8\Z$, the distinct points $0$ and $2$ have difference $2$,
which has no inverse. Indeed, $f(2)-f(0)$ is even for every polynomial,
so the data $f(0)=0$, $f(2)=1$ cannot be interpolated.
\end{warningbox}

\begin{checkpoint}
Both reconstruction formulas first make an object that vanishes at every
unwanted coordinate, then multiply by an inverse to make its remaining
coordinate equal to $1$. Identify those two steps before memorizing either
formula.
\end{checkpoint}
